\documentclass[b4paper,7pt,twocolumn,landscape]{jsarticle}
\usepackage{info_literacy_exam}

\begin{document}

\twocolumn[
\begin{center}
\large{\textbf{情報リテラシー 2025年度 学年末試験 問題用紙}}
\end{center}

\noindent
出席番号：\underline{\hspace{2cm}} \quad 氏名：\underline{\hspace{6cm}}

\vspace{0.5mm}
]

\subsubsection*{注意事項}
{
\begin{itemize}
  \item マークシートの学年は、\textbf{1}を記入・マークしてください。
  \item マークシートのクラスは、M:01、E:02、I:03、C:04、A:05の該当する数字を記入・マークしてください。
  \item マークシートの番号は、出席番号を\textbf{3桁}に変換して、右詰めにして、記入・マークしてください。
  \item 問題は全部で\textbf{50問}あります。\textbf{100点満点}（各2点）です。
  \item 解答は\textbf{マークシート}の解答欄に\textbf{正確に濃く記入}してください。
  \item すべて4択問題です。各設問について最も適切な選択肢を1つ選び、マークしてください。
  \item 到達目標5：問1$\sim$問50
\end{itemize}
}

\subsubsection*{【第8回】データサイエンス・AI：概要（問1〜問6）}

\begin{enumerate}
\item データサイエンスの定義として最も適切なものはどれか。\\
\choices{ハードウェアを設計する技術|データから価値ある情報を引き出す技術|セキュリティを守る技術|Webサイトをデザインする技術}

\item AI（人工知能）の定義として最も適切なものはどれか。\\
\choices{データを保存する技術|人間の知的活動をコンピュータで実現する技術|インターネットに接続する技術|電気回路を設計する技術}

\item Google Colabの説明として\textbf{誤っている}ものはどれか。\\
\choices{ブラウザ上でPythonを実行できる|Googleアカウントがあれば無料で使える|専用ソフトウェアのインストールが必要である|クラウドで動くのでPCのスペックに依存しない}

\item データサイエンスでグラフを描く理由として最も適切なものはどれか。\\
\choices{プログラミングの練習のため|傾向や特徴を一瞬で理解できるため|ファイルサイズを小さくするため|計算速度を向上させるため}

\item 次のうち、身近なAI活用例として\textbf{誤っている}ものはどれか。\\
\choices{スマホの顔認証|YouTubeのおすすめ動画|電卓の四則演算|スパムメールフィルター}

\item AI倫理が重要とされる理由として最も適切なものはどれか。\\
\choices{計算速度を上げるため|間違った使い方で人が傷つく可能性があるため|プログラムを短くするため|消費電力を減らすため}
\end{enumerate}

\subsubsection*{【第9回】AI活用事例とその倫理的課題（問7〜問13）}

\begin{enumerate}
\setcounter{enumi}{6}
\item AIにおける「バイアス」の意味として最も適切なものはどれか。\\
\choices{AIの計算速度|データや判断の偏り|AIの消費電力|プログラムの長さ}

\item ある企業の採用AIが女性を低評価した原因として最も適切なものはどれか。\\
\choices{ハードウェアが故障していた|過去データが男性中心で「男性＝良い候補者」と学習|女性の応募者数が多すぎた|プログラムにバグがあった}

\item 「データの偏り＝AIの偏り」という説明として最も適切なものはどれか。\\
\choices{AIはデータとは無関係に判断する|データに偏りがあるとAIの判断も偏る|AIは常に公平な判断をする|データが多ければバイアスは発生しない}

\item 顔認識技術の活用例として\textbf{誤っている}ものはどれか。\\
\choices{空港の入国審査|スマホのロック解除|天気予報の精度向上|商業施設の防犯}

\item 顔認識カメラによる監視社会のリスクとして最も適切なものはどれか。\\
\choices{計算速度が遅くなる|プライバシー侵害と行動の自由の制限|電力消費が増える|インターネット接続が不安定になる}

\item AIの「ブラックボックス問題」の説明として最も適切なものはどれか。\\
\choices{AIの筐体が黒い色をしていること|AIがなぜその判断をしたかが分からないこと|AIが暗い場所で動作しないこと|AIのデータが暗号化されていること}

\item 医療診断AIでブラックボックス問題が特に問題となる理由は何か。\\
\choices{診断理由を医師が患者に説明できないため|AIの計算が遅くなるため|病院の電気代が増えるため|患者の名前が分からなくなるため}
\end{enumerate}

\subsubsection*{【第10回】データ分析と仮説検証サイクル（問14〜問21）}

\begin{enumerate}
\setcounter{enumi}{13}
\item 仮説検証サイクルの4ステップの正しい順序はどれか。\\
\choices{結論→検証→仮説→分析|仮説→検証→分析→結論|分析→仮説→結論→検証|検証→結論→分析→仮説}

\item 相関係数の取りうる範囲として正しいものはどれか。\\
\choices{0〜100|0〜1|$-1$〜$+1$|$-100$〜$+100$}

\item 散布図の説明として最も適切なものはどれか。\\
\choices{時間による変化を見るグラフ|2つのデータの関係を見るグラフ|割合の違いを把握するグラフ|数値の大小を比較するグラフ}

\item 「相関＝因果ではない」の意味として最も適切なものはどれか。\\
\choices{相関が高ければ必ず因果関係がある|相関が高くても因果とは断定できない|相関と因果は全く同じ意味である|因果関係があれば相関はゼロになる}

\item 線形回帰の式 $y = ax + b$ において、$a$ は何を表すか。\\
\choices{切片|傾き|相関係数|決定係数}

\item 決定係数（$R^2$）が0.6のとき、その解釈として最も適切なものはどれか。\\
\choices{モデルがデータの約60\%を説明できている|モデルの精度が60点である|データの60\%が外れ値である|相関係数が0.6である}

\item データを訓練データとテストデータに分ける理由として最も適切なものはどれか。\\
\choices{データ量を減らすため|未知のデータでも予測できるか確認するため|計算を速くするため|データを暗号化するため}

\item 外れ値がある場合の影響として\textbf{誤っている}ものはどれか。\\
\choices{回帰直線がずれる|予測精度が下がる|必ずデータから削除すべきである|新しい発見につながる可能性がある}
\end{enumerate}

\subsubsection*{【第11回】機械学習の基礎（問22〜問29）}

\begin{enumerate}
\setcounter{enumi}{21}
\item 機械学習の定義として最も適切なものはどれか。\\
\choices{機械を物理的に修理する技術|データから学習し予測や判断を行う技術|機械の電源を入れる技術|機械を製造する技術}

\newpage

\item 機械学習の3つの分類として正しい組み合わせはどれか。\\
\choices{教師あり学習、教師なし学習、強化学習|入力学習、出力学習、中間学習|高速学習、低速学習、中速学習|画像学習、音声学習、テキスト学習}

\item 教師あり学習の特徴として最も適切なものはどれか。\\
\choices{正解ラベルなしでパターンを発見する|正解ラベル付きデータで学習し予測する|試行錯誤で最適な行動を学習する|データを使わずに学習する}

\item 教師あり学習における「回帰」と「分類」の違いとして正しいものはどれか。\\
\choices{回帰は画像を扱い、分類はテキストを扱う|回帰は数値を予測し、分類はカテゴリを予測する|回帰は高速で、分類は低速である|回帰と分類に違いはない}

\item 次のうち、教師あり学習の「分類」の例として最も適切なものはどれか。\\
\choices{気温から売上（円）を予測|メールがスパムか正常かを判定|似た顧客をグループ化|ゲームで高得点を目指す}

\item 教師なし学習の「クラスタリング」の説明として最も適切なものはどれか。\\
\choices{正解ラベルを使って分類する|正解ラベルなしで似たものをグループ化する|報酬を最大化する行動を学習する|数値を予測する}

\item 強化学習の具体例として最も適切なものはどれか。\\
\choices{スパムメール判定|顧客セグメンテーション|囲碁AI（AlphaGo）が試行錯誤で強くなる|手書き数字の認識}

\item K-means法で変数選択が重要な理由として最も適切なものはどれか。\\
\choices{計算時間を短くするため|互いに異なる情報を持つ変数を選ぶことでグループ分けがうまくいくため|変数を増やすほど精度が上がるため|変数選択は結果に影響しないため}
\end{enumerate}

\subsubsection*{【第12回】深層学習（問30〜問37）}

\begin{enumerate}
\setcounter{enumi}{29}
\item 深層学習（ディープラーニング）の定義として最も適切なものはどれか。\\
\choices{データを深く保存する技術|ニューラルネットワークを多層化した機械学習|深い場所でコンピュータを動かす技術|プログラムを深く理解する技術}

\item MNISTデータセットの訓練データの枚数として正しいものはどれか。\\
\choices{6枚|60枚|60,000枚|600,000枚}

\item 従来の機械学習と深層学習の違いとして最も適切なものはどれか。\\
\choices{従来は特徴量を人間が設計、深層学習は自動で学習|深層学習は特徴量を人間が設計、従来は自動で学習|両者に違いはない|従来の機械学習の方が大量データを必要とする}

\item ニューラルネットワークの基本構造として正しい順序はどれか。\\
\choices{出力層→隠れ層→入力層|隠れ層→入力層→出力層|入力層→隠れ層→出力層|入力層→出力層→隠れ層}

\item 活性化関数ReLU（レルー）の特徴として最も適切なものはどれか。\\
\choices{出力を確率に変換する|マイナスの値を0にする|常に1を出力する|入力をそのまま出力する}

\item 活性化関数Softmax（ソフトマックス）の特徴として最も適切なものはどれか。\\
\choices{マイナスの値を0にする|出力を確率に変換する（合計が1になる）|入力を2倍にする|常に0を出力する}

\item MNISTの画像サイズが28×28ピクセルのとき、入力層のノード数は何か。\\
\choices{28|56|784|280}

\item 「エポック」の意味として最も適切なものはどれか。\\
\choices{テストデータの数|訓練データ全体を1回学習すること|ニューラルネットワークの層の数|出力層のノード数}
\end{enumerate}

\subsubsection*{【第13回】生成AIと倫理的課題（問38〜問43）}

\begin{enumerate}
\setcounter{enumi}{37}
\item 生成AIの定義として最も適切なものはどれか。\\
\choices{データを削除するAI|新しいコンテンツ（テキスト、画像など）を生成するAI|コンピュータを修理するAI|ネットワークを監視するAI}

\item 生成AIの「ハルシネーション」の意味として最も適切なものはどれか。\\
\choices{AIが高速に動作すること|事実でない情報をもっともらしく生成すること|AIが画像を生成すること|AIが音声を認識すること}

\item 生成AIの倫理的課題として\textbf{誤っている}ものはどれか。\\
\choices{バイアスと公平性の問題|著作権の問題|計算速度が速すぎる問題|プライバシーの問題}

\item NotebookLMとCopilotの違いとして最も適切なものはどれか。\\
\choices{NotebookLMは画像生成専用、Copilotはテキスト生成専用|NotebookLMはアップロード資料のみを参照、Copilotは一般知識を参照|NotebookLMは無料、Copilotは有料のみ|両者に違いはない}

\item 生成AIを使う際の注意点として最も適切なものはどれか。\\
\choices{AIの回答は常に正しいのでそのまま信じてよい|AIの回答を鵜呑みにせず批判的に読み最終判断は自分で行う|AIに機密情報を積極的に入力すべき|AIを使えば自分で考える必要はない}

\item 次のうち、テキスト生成AIの例として\textbf{誤っている}ものはどれか。\\
\choices{ChatGPT|Claude|Copilot|Stable Diffusion}
\end{enumerate}

\subsubsection*{【第14回】AIモデルの評価（問44〜問50）}

\begin{enumerate}
\setcounter{enumi}{43}
\item 「推論」の定義として最も適切なものはどれか。\\
\choices{モデルを学習させること|学習したモデルで新しいデータを予測すること|データを収集すること|プログラムを書くこと}

\item 正解率（Accuracy）の計算式として正しいものはどれか。\\
\choices{正解数÷予測した数|正解数÷全体の数|予測した数÷正解数|全体の数÷正解数}

\item 適合率（Precision）の意味として最も適切なものはどれか。\\
\choices{実際のそのクラスをどれだけ見つけられたか|「このクラスだ」と予測したとき実際に正しい割合|全体で何\%正解したか|モデルの学習速度}

\item 再現率（Recall）の意味として最も適切なものはどれか。\\
\choices{「このクラスだ」と予測したとき実際に正しい割合|実際のそのクラスをどれだけ見つけられたか（見逃しの少なさ）|全体で何\%正解したか|モデルのパラメータ数}

\item 「7だ」と予測1,050回、うち実際に7だったのが1,000回のとき、適合率は約何\%か。\\
\choices{約90.5\%|約95.2\%|約97.3\%|約100\%}

\item 過学習の特徴として最も適切なものはどれか。\\
\choices{訓練データでは精度が高いがテストデータでは低い|訓練データでもテストデータでも精度が低い|訓練データでもテストデータでも精度が高い|学習に時間がかからない}

\item 病気診断AIにおいて「再現率」を重視すべき理由として最も適切なものはどれか。\\
\choices{誤診を減らすため|計算速度を上げるため|消費電力を減らすため|病気の人を見逃さないため}
\end{enumerate}

\end{document}
